\documentclass[10pt,a4paper]{amsart}
\usepackage{amsmath,amssymb,amsthm,graphicx}
\usepackage{verbatim}
\newenvironment{reference}{\comment}{\endcomment}
%\newenvironment{reference}{}{}
\newenvironment{slogan}{\comment}{\endcomment}
\newenvironment{history}{\comment}{\endcomment}

% For commutative diagrams we use Xy-pic
\usepackage[all]{xy}

% We use 2cell for 2-commutative diagrams.
\xyoption{2cell}
% Theorem environments.
%
\theoremstyle{plain}
\newtheorem{theorem}[subsection]{Theorem}
\newtheorem{proposition}[subsection]{Proposition}
\newtheorem{lemma}[subsection]{Lemma}

\theoremstyle{definition}
\newtheorem{definition}[subsection]{Definition}
\newtheorem{example}[subsection]{Example}
\newtheorem{exercise}[subsection]{Exercise}
\newtheorem{situation}[subsection]{Situation}

\theoremstyle{remark}
\newtheorem{remark}[subsection]{Remark}
\newtheorem{remarks}[subsection]{Remarks}

\numberwithin{equation}{subsection}

% Macros
%
\def\lim{\mathop{\mathrm{lim}}\nolimits}
\def\colim{\mathop{\mathrm{colim}}\nolimits}
\def\Spec{\mathop{\mathrm{Spec}}}
\def\Hom{\mathop{\mathrm{Hom}}\nolimits}
\def\Ext{\mathop{\mathrm{Ext}}\nolimits}
\def\SheafHom{\mathop{\mathcal{H}\!\mathit{om}}\nolimits}
\def\SheafExt{\mathop{\mathcal{E}\!\mathit{xt}}\nolimits}
\def\Sch{\mathit{Sch}}
\def\Mor{\mathop{\mathrm{Mor}}\nolimits}
\def\Ob{\mathop{\mathrm{Ob}}\nolimits}
\def\Sh{\mathop{\mathit{Sh}}\nolimits}
\def\NL{\mathop{N\!L}\nolimits}
\def\CH{\mathop{\mathrm{CH}}\nolimits}
\def\proetale{{pro\text{-}\acute{e}tale}}
\def\etale{{\acute{e}tale}}
\def\QCoh{\mathit{QCoh}}
\def\Ker{\mathop{\mathrm{Ker}}}
\def\Im{\mathop{\mathrm{Im}}}
\def\Coker{\mathop{\mathrm{Coker}}}
\def\Coim{\mathop{\mathrm{Coim}}}

% Boxtimes
%
\DeclareMathSymbol{\boxtimes}{\mathbin}{AMSa}{"02}

%
% Macros for moduli stacks/spaces
%
\def\QCohstack{\mathcal{QC}\!\mathit{oh}}
\def\Cohstack{\mathcal{C}\!\mathit{oh}}
\def\Spacesstack{\mathcal{S}\!\mathit{paces}}
\def\Quotfunctor{\mathrm{Quot}}
\def\Hilbfunctor{\mathrm{Hilb}}
\def\Curvesstack{\mathcal{C}\!\mathit{urves}}
\def\Polarizedstack{\mathcal{P}\!\mathit{olarized}}
\def\Complexesstack{\mathcal{C}\!\mathit{omplexes}}
% \Pic is the operator that assigns to X its picard group, usage \Pic(X)
% \Picardstack_{X/B} denotes the Picard stack of X over B
% \Picardfunctor_{X/B} denotes the Picard functor of X over B
\def\Pic{\mathop{\mathrm{Pic}}\nolimits}
\def\Picardstack{\mathcal{P}\!\mathit{ic}}
\def\Picardfunctor{\mathrm{Pic}}
\def\Deformationcategory{\mathcal{D}\!\mathit{ef}}
\usepackage{cleveref}
\crefname{originallemma}{original lemma}{original lemmas}
\Crefname{originallemma}{Original lemma}{Original lemmas}
\makeatletter
\newlabel{sites-lemma-colimit-over-ordinal-sections}{{0GS0}{source}{}{sourcepointer-1}{}}
\newlabel{sites-lemma-colimit-over-ordinal-sections@cref}{{[originallemma][0][]0GS0}{[1][0][]source}{}{}{}}
\makeatother
\crefname{originallemma}{original lemma}{original lemmas}
\Crefname{originallemma}{Original lemma}{Original lemmas}
\makeatletter
\newlabel{sets-lemma-map-from-set-lifts}{{05N2}{source}{}{sourcepointer-2}{}}
\newlabel{sets-lemma-map-from-set-lifts@cref}{{[originallemma][0][]05N2}{[1][0][]source}{}{}{}}
\makeatother
\crefname{originalproposition}{original proposition}{original propositions}
\Crefname{originalproposition}{Original proposition}{Original propositions}
\makeatletter
\newlabel{sets-proposition-exist-ordinals-large-cofinality}{{05N3}{source}{}{sourcepointer-3}{}}
\newlabel{sets-proposition-exist-ordinals-large-cofinality@cref}{{[originalproposition][0][]05N3}{[1][0][]source}{}{}{}}
\makeatother
\title{Every map from a fixed family of sheaves factors through an earlier stage}
\author{The Stacks Project Authors}
\date{stacks-project:injectives.tex:local-snapshot-20261012; GNU FDL 1.2 or later}
\begin{document}
\maketitle
\noindent\textit{Editorial note.} All delivered mathematical text is the author's own, copied from the pinned Stacks Project chapter file sites-cohomology.tex; no statement or proof step was rewritten, reordered or removed.\par\medskip
\noindent\textit{Editorial note.} The proof is the author's complete argument; the counting step and the two cited inputs are all named, and no step is left to the reader.\par\medskip
\noindent\textit{Editorial note.} Every labelled reference occurring inside the counted spans is carried in assembly.external\_references as a non-load-bearing navigation pointer (same-chapter ones at their definition site, cross-chapter ones at the invocation inside the counted span), so no load-bearing mathematics is replaced by a stub. No bibliography citation occurs inside the counted spans.\par\medskip
\noindent\textit{Editorial note.} The same-name check for this card scanned every staged Stacks chapter of this repository -- all sources/stacks-*/pinned-original-source.tex files, the card's own chapter excluded -- and the labels found are recorded in dedup.same\_name\_cross\_chapter on the card.\par\medskip
\noindent\textit{Editorial note.} The preamble is the chapter's own retained minimal preamble (sources/stacks-sites-cohomology/retained-preamble.tex), a verbatim excerpt of the native preamble: the Xy-pic setup and the native theorem declarations with the mathematical macros.\par\medskip
\noindent\textit{Allowed standard prerequisites (editorial).}\begin{itemize}
\item The transfinite construction of the injective replacements. Same chapter, the construction of the sheaves attached to an ordinal and the functor used in it. Supplies the stages through which the factorisation is asserted and the inclusions between them.
\item Cardinals of maps out of a set. Sets, Lemma sets-lemma-map-from-set-lifts. Supplies the lifting statement that bounds the set of maps needing to be killed at each stage.
\item Ordinals of arbitrarily large cofinality. Sets, Proposition sets-proposition-exist-ordinals-large-cofinality. Supplies the existence of an ordinal of cofinality larger than the relevant cardinal, which makes the bound effective.
\end{itemize}\medskip
\ifdefined\hypertarget\hypertarget{sourcepointer-1}{}\fi
\noindent\textit{Source reference.} Original lemma 0GS0: sources/stacks-sites/pinned-original-source.tex, source lines 3921--3921. Navigation pointer invoked at proof line 984; the counted span contains the pointer only and the cited result is not reproduced inside it, so the pointer is registered as a non-load-bearing navigation stub in the established form for Stacks material. No bibliography citation is involved.\par\medskip
\ifdefined\hypertarget\hypertarget{sourcepointer-2}{}\fi
\noindent\textit{Source reference.} Original lemma 05N2: sources/stacks-sets/pinned-original-source.tex, source lines 410--410. Navigation pointer invoked at proof line 989; the counted span contains the pointer only and the cited result is not reproduced inside it, so the pointer is registered as a non-load-bearing navigation stub in the established form for Stacks material. No bibliography citation is involved.\par\medskip
\ifdefined\hypertarget\hypertarget{sourcepointer-3}{}\fi
\noindent\textit{Source reference.} Original proposition 05N3: sources/stacks-sets/pinned-original-source.tex, source lines 434--434. Navigation pointer invoked at proof line 994; the counted span contains the pointer only and the cited result is not reproduced inside it, so the pointer is registered as a non-load-bearing navigation stub in the established form for Stacks material. No bibliography citation is involved.\par\medskip
\begin{lemma}
\label{lemma-map-into-smaller}
Suppose that $\mathcal{G}_i$, $i\in I$ is set of abelian sheaves
on $\mathcal{C}$. There exists an ordinal $\beta$ such that
for any sheaf $\mathcal{F}$, any $i\in I$, and any map
$\varphi : \mathcal{G}_i \to J_\beta(\mathcal{F})$ there exists an
$\alpha < \beta$ such that $ \varphi $ factors through
$J_\alpha(\mathcal{F})$.
\end{lemma}

\begin{proof}
This reduces to the case of a single sheaf $\mathcal{G}$
by taking the direct sum of all the $\mathcal{G}_i$.

\medskip\noindent
Consider the sets
$$
S = \coprod\nolimits_{U \in \Ob(\mathcal{C})} \mathcal{G}(U).
$$
and
$$
T_\beta
=
\coprod\nolimits_{U \in \Ob(\mathcal{C})} J_\beta(\mathcal{F})(U)
$$
The transition maps between the sets $T_\beta$ are injective.
If the cofinality of $\beta$ is large enough, then
$T_\beta = \colim_{\alpha < \beta} T_\alpha$, see
Sites, Lemma \ref{sites-lemma-colimit-over-ordinal-sections}.
A morphism $\mathcal{G} \to J_\beta(\mathcal{F})$ factors
through $J_\alpha(\mathcal{F})$ if and only if
the associated map $S \to T_\beta$ factors through $T_\alpha$.
By
Sets, Lemma \ref{sets-lemma-map-from-set-lifts}
if the cofinality of $\beta$ is bigger than the cardinality
of $S$, then the result of the lemma is true. Hence the lemma
follows from the fact that there are ordinals with arbitrarily
large cofinality, see
Sets, Proposition \ref{sets-proposition-exist-ordinals-large-cofinality}.
\end{proof}

\end{document}
